% GENERATED FILE. Edit canonical lessons, then run npm run generate:books.
% canonical-source-hash: fnv1a64:631a04dcd79eddf7
% canonical-lessons: PT-C422-porao, PT-C422-balanca, PT-C422-empacotar, PT-C422-cadeado, PT-C422-extraviar

\chapter{As bagagens}
\label{ch:pt-as-bagagens}
\hlchaptermodality{422}

\begin{chapteropening}
\textbf{By the end of this chapter:} \emph{I can deal with luggage: the hold, the scales, packing up, the padlock and bags that go missing.}

Complete the last lesson of chapter 422: I can deal with luggage: the hold, the scales, packing up, the padlock and bags that go missing.
\end{chapteropening}

\section[o porão]{o porão — the hold (of a plane or ship); the cellar}

\label{lesson:PT-C422-porao}



\begin{quote}
Before the new one: say the Portuguese for the runway, then the Portuguese for the time zone.
\end{quote}

\subsection*{You'll want to know: o porão}
\textbf{o porão} — \textquotedblleft{}the hold (of a plane or ship); the cellar\textquotedblright{}.

\begin{quote}
A mala grande vai no porão.
\end{quote}

\begin{quote}
\emph{The big suitcase goes in the hold.}
\end{quote}

In Brazil \emph{o porão} is mostly a basement.

\begin{cousinweb}
\emph{porão} comes from Old Portuguese \emph{prão} (EtymDB); its earlier origin is not settled.
\end{cousinweb}

\begin{grammarlens}[title={What you've built}]
One more word for telling, explaining and describing at B1.
\end{grammarlens}

\subsection*{Your turn}
\begin{itemize}
\raggedright
  \item \emph{Say it:} \emph{o porão}
  \item \emph{Say it:} \emph{o porão}, once more
  \item \emph{Say it:} a sentence of your own with \emph{o porão}
  \item \emph{Recall:} say the Portuguese for the runway, then the Portuguese for the time zone, then say \emph{o porão} again
\end{itemize}

\begin{tcolorbox}[breakable,colback=teal!4,colframe=teal!35!black,title={Before you move on}]
What does o porão mean? (\textquotedblleft{}The hold (of a plane or ship); the cellar\textquotedblright{}.) Say it once more.
\end{tcolorbox}
\section[a balança]{a balança — the scales}

\label{lesson:PT-C422-balanca}



\begin{quote}
Before the new one: say the Portuguese for the time zone, then the Portuguese for the hold.
\end{quote}

\subsection*{You'll want to know: a balança}
\textbf{a balança} — \textquotedblleft{}the scales\textquotedblright{}.

\begin{quote}
No aeroporto, a balança mostra vinte e três quilos.
\end{quote}

\begin{quote}
\emph{At the airport, the scales show twenty-three kilos.}
\end{quote}

\emph{A balança comercial} — the balance of trade.

\begin{cousinweb}
\emph{balança} is inherited from Old Portuguese \emph{balança} (EtymDB), from Vulgar Latin \emph{bilancia}, built on Latin \emph{bilanx} 'having two scale pans' (Lewis \& Short).
\end{cousinweb}

\begin{grammarlens}[title={What you've built}]
One more word for telling, explaining and describing at B1.
\end{grammarlens}

\subsection*{Your turn}
\begin{itemize}
\raggedright
  \item \emph{Say it:} \emph{a balança}
  \item \emph{Say it:} \emph{a balança}, once more
  \item \emph{Say it:} a sentence of your own with \emph{a balança}
  \item \emph{Recall:} say the Portuguese for the time zone, then the Portuguese for the hold, then say \emph{a balança} again
\end{itemize}

\begin{tcolorbox}[breakable,colback=teal!4,colframe=teal!35!black,title={Before you move on}]
What does a balança mean? (\textquotedblleft{}The scales\textquotedblright{}.) Say it once more.
\end{tcolorbox}
\section[empacotar]{empacotar — to pack up, to wrap into parcels}

\label{lesson:PT-C422-empacotar}



\begin{quote}
Before the new one: say the Portuguese for the hold, then the Portuguese for the scales.
\end{quote}

\subsection*{You'll want to know: empacotar}
\textbf{empacotar} — \textquotedblleft{}to pack up, to wrap into parcels\textquotedblright{}.

\begin{quote}
Empacotamos os livros em caixas de papel.
\end{quote}

\begin{quote}
\emph{We pack up the books in cardboard boxes.}
\end{quote}

For a suitcase people usually say \emph{fazer a mala}.

\begin{cousinweb}
\emph{empacotar} is built from \emph{em-} and \emph{pacote} 'packet', a small form of \emph{paca} 'bundle' (EtymDB).
\end{cousinweb}

\begin{grammarlens}[title={What you've built}]
One more word for telling, explaining and describing at B1.
\end{grammarlens}

\subsection*{Your turn}
\begin{itemize}
\raggedright
  \item \emph{Say it:} \emph{empacotar}
  \item \emph{Say it:} \emph{empacotar}, once more
  \item \emph{Say it:} a sentence of your own with \emph{empacotar}
  \item \emph{Recall:} say the Portuguese for the hold, then the Portuguese for the scales, then say \emph{empacotar} again
\end{itemize}

\begin{tcolorbox}[breakable,colback=teal!4,colframe=teal!35!black,title={Before you move on}]
What does empacotar mean? (\textquotedblleft{}To pack up, to wrap into parcels\textquotedblright{}.) Say it once more.
\end{tcolorbox}
\section[o cadeado]{o cadeado — the padlock}

\label{lesson:PT-C422-cadeado}



\begin{quote}
Before the new one: say the Portuguese for the scales, then the Portuguese for to pack up.
\end{quote}

\subsection*{You'll want to know: o cadeado}
\textbf{o cadeado} — \textquotedblleft{}the padlock\textquotedblright{}.

\begin{quote}
Fechei a mochila com um cadeado.
\end{quote}

\begin{quote}
\emph{I locked my backpack with a padlock.}
\end{quote}

\emph{O cadeado da bicicleta} — the bike lock.

\begin{cousinweb}
\emph{cadeado} is inherited from Late Latin \emph{catēnātum} 'chained' (EtymDB), built on \emph{catēna} 'chain'.
\end{cousinweb}

\begin{grammarlens}[title={What you've built}]
One more word for telling, explaining and describing at B1.
\end{grammarlens}

\subsection*{Your turn}
\begin{itemize}
\raggedright
  \item \emph{Say it:} \emph{o cadeado}
  \item \emph{Say it:} \emph{o cadeado}, once more
  \item \emph{Say it:} a sentence of your own with \emph{o cadeado}
  \item \emph{Recall:} say the Portuguese for the scales, then the Portuguese for to pack up, then say \emph{o cadeado} again
\end{itemize}

\begin{tcolorbox}[breakable,colback=teal!4,colframe=teal!35!black,title={Before you move on}]
What does o cadeado mean? (\textquotedblleft{}The padlock\textquotedblright{}.) Say it once more.
\end{tcolorbox}
\section[extraviar]{extraviar — to lose, to mislay (luggage or post); to go astray}

\label{lesson:PT-C422-extraviar}



\begin{quote}
Before the new one: say the Portuguese for to pack up, then the Portuguese for the padlock.
\end{quote}

\subsection*{You'll want to know: extraviar}
\textbf{extraviar} — \textquotedblleft{}to lose, to mislay (luggage or post); to go astray\textquotedblright{}.

\begin{quote}
A companhia aérea extraviou a minha mala.
\end{quote}

\begin{quote}
\emph{The airline lost my suitcase.}
\end{quote}

\emph{A bagagem extraviada} — lost luggage.

\begin{cousinweb}
\emph{extraviar} is built from \emph{extra-} 'outside' and \emph{via} 'way', from Latin \emph{via} (Lewis \& Short): to send off the right road.
\end{cousinweb}

\begin{grammarlens}[title={What you've built}]
One more word for telling, explaining and describing at B1.
\end{grammarlens}

\subsection*{Your turn}
\begin{itemize}
\raggedright
  \item \emph{Say it:} \emph{extraviar}
  \item \emph{Say it:} \emph{extraviar}, once more
  \item \emph{Say it:} a sentence of your own with \emph{extraviar}
  \item \emph{Recall:} say the Portuguese for to pack up, then the Portuguese for the padlock, then say \emph{extraviar} again
\end{itemize}

\begin{tcolorbox}[breakable,colback=teal!4,colframe=teal!35!black,title={Before you move on}]
What does extraviar mean? (\textquotedblleft{}To lose, to mislay (luggage or post); to go astray\textquotedblright{}.) Say it once more.
\end{tcolorbox}
